\section{Related work}\label{sec:literature-review-related-work}

For additional readings, the reader is referred to studies related to ours, as detailed below.

\textbf{Closest to our work} is the survey by \textcite{olsson2023systems}, which reviews ten studies combining DTs and SoS and identifies challenges related to interoperability and lifecycle management. Our study expands this scope through a systematic analysis of 80 primary studies, examining motivations, architectures, technical characteristics, non-functional properties, and research maturity to clarify how DT architectures support SoS characteristics in SoTS.

Several studies acknowledge that DTs can be composed into SoS and that such integrated systems may bring unique benefits. \textcite{mylonas2021digital} survey DTs in manufacturing and smart cities, and investigate their integration patterns. \textcite{semeraro2021digital} conduct a cross-domain systematic literature review of DTs and introduce the patterns of digital twinning at the unit and system level of SoS. \textcite{bottjer2023review} systematically review unit level manufacturing DTs, situating them within hierarchical system and SoS structures and identifying interoperability and standardization as prerequisites for large-scale integration. \textcite{bertoni2022designing} survey DTs in product-service systems and argue that analyzing emergent behavior and lifecycle value requires connecting multiple DTs into higher-order SoS configurations. \textcite{dalibor2022cross} present a systematic mapping study of DT research and identify systems of systems as a class of physical systems that can be digitally twinned. \textcite{klar2023digital} review existing DT maturity models and propose a framework in which interoperability across DTs represents the highest level of maturity, enabling operation as an SoS.

Despite the benefits, engineering DTs of SoS remains challenged by interoperability gaps, connectivity and privacy concerns, limited standardization, and the operational independence of constituent systems \cite{michael2022integration,dietz2020digital,borth2019digital}. These challenges make the composition of DT functions and behaviors across interconnected systems difficult in practice \cite{gill2024toward}.